\addcontentsline{toc}{section}{СПИСОК ИСПОЛЬЗОВАННЫХ ИСТОЧНИКОВ}

% Список оформляется по ГОСТ Р 7.0.100-2018. Источники должны быть преимущественно
% не старше 2014 года. Во второй итерации необходимо проверить и при необходимости
% обновить даты обращения к электронным ресурсам.

\begin{thebibliography}{99}

% --- ГОСТы ---
\bibitem{gost19102} ГОСТ 19.102--77. Единая система программной документации. Стадии разработки~: межгосударственный стандарт~: издание официальное~: утвержден Постановлением Государственного комитета стандартов Совета Министров СССР от 20 декабря 1977~г. №~2851~: дата введения 1980-01-01.~--~Москва~: Стандартинформ, 2010.~--~Текст~: непосредственный.

\bibitem{gost34601} ГОСТ 34.601--90. Информационная технология. Комплекс стандартов на автоматизированные системы. Автоматизированные системы. Стадии создания~: межгосударственный стандарт~: издание официальное~: утвержден и введен в действие Постановлением Государственного комитета СССР по управлению качеством продукции и стандартам от 29 декабря 1990~г. №~3469~: дата введения 1992-01-01.~--~Москва~: Стандартинформ, 2009.~--~Текст~: непосредственный.

\bibitem{gost7320} ГОСТ 7.32--2017. Система стандартов по информации, библиотечному и издательскому делу. Отчет о научно-исследовательской работе. Структура и правила оформления~: межгосударственный стандарт~: издание официальное~: введен в действие 2018-07-01.~--~Москва~: Стандартинформ, 2017.~--~Текст~: непосредственный.

\bibitem{gost70100} ГОСТ Р 7.0.100--2018. Система стандартов по информации, библиотечному и издательскому делу. Библиографическая запись. Библиографическое описание. Общие требования и правила составления~: национальный стандарт Российской Федерации~: издание официальное~: введен в действие 2019-07-01.~--~Москва~: Стандартинформ, 2018.~--~Текст~: непосредственный.

\bibitem{gost25010} ГОСТ Р ИСО/МЭК 25010--2015. Информационные технологии. Системная и программная инженерия. Требования и оценка качества систем и программного обеспечения (SQuaRE). Модели качества систем и программных продуктов~: национальный стандарт Российской Федерации~: издание официальное~: утвержден и введен в действие Приказом Федерального агентства по техническому регулированию и метрологии от 29 мая 2015~г. №~442-ст~: дата введения 2016-06-01.~--~Москва~: Стандартинформ, 2015.~--~Текст~: непосредственный.

\bibitem{stu02030} СТУ 02.030--2023. Курсовые работы (проекты). Выпускные квалификационные работы. Общие требования к структуре и оформлению~: стандарт университета.~--~Курск~: Юго-Западный государственный университет, 2023.~--~Текст~: непосредственный.

% --- UML, REST, веб-программирование ---
\bibitem{omguml} OMG Unified Modeling Language (OMG UML). Version 2.5.1~: specification / Object Management Group.~--~Needham~: OMG, 2017.~--~796~p.~--~URL: \url{https://www.omg.org/spec/UML/2.5.1} (дата обращения: 10.05.2026).~--~Текст~: электронный.

\bibitem{rfc9110} Hypertext Transfer Protocol (HTTP) Semantics~: RFC~9110 / R. Fielding, M. Nottingham, J. Reschke.~--~Internet Engineering Task Force, June 2022.~--~URL: \url{https://datatracker.ietf.org/doc/html/rfc9110} (дата обращения: 10.05.2026).~--~Текст~: электронный.

\bibitem{rfc6455} The WebSocket Protocol~: RFC~6455 / I. Fette, A. Melnikov.~--~Internet Engineering Task Force, December 2011.~--~URL: \url{https://datatracker.ietf.org/doc/html/rfc6455} (дата обращения: 10.05.2026).~--~Текст~: электронный.

\bibitem{owasp} OWASP Top 10:2021 // The OWASP Foundation~: [сайт].~--~2021.~--~URL: \url{https://owasp.org/Top10/2021/} (дата обращения: 10.05.2026).~--~Текст~: электронный.

\bibitem{kovalev} Ковалев, В.~В. Актуальность использования архитектуры REST для обмена данными между клиент-серверными приложениями / В.~В.~Ковалев, В.~В.~Храмков, Е.~И.~Семенова.~--~Текст~: электронный // Актуальные проблемы авиации и космонавтики.~--~2019.~--~Т.~2.~--~URL: \url{https://cyberleninka.ru/article/n/aktualnost-ispolzovaniya-arhitektury-rest-dlya-obmena-dannymi-mezhdu-klient-servernymi-prilozheniyami} (дата обращения: 10.05.2026).

\bibitem{date} Дейт, К.~Дж. Введение в системы баз данных / К.~Дж.~Дейт~; перевод с английского.~--~8-е изд.~--~Москва~: ООО <<И.~Д.~Вильямс>>, 2005.~--~1328~с.~--~ISBN 5-8459-0788-8.~--~Текст~: непосредственный.

\bibitem{connolly} Коннолли, Т. Базы данных. Проектирование, реализация и сопровождение. Теория и практика / Т.~Коннолли, К.~Бегг~; перевод с английского.~--~3-е изд.~--~Москва~: ООО <<И.~Д.~Вильямс>>, 2017.~--~1440~с.~--~ISBN 978-5-8459-2020-1.~--~Текст~: непосредственный.

\bibitem{mysqldoc} MySQL 8.0 Reference Manual~: официальный сайт.~--~Oracle Corporation, 1995--.~--~URL: \url{https://dev.mysql.com/doc/refman/8.0/en/} (дата обращения: 10.05.2026).~--~Текст~: электронный.

\bibitem{romanov} Романов, С.~С. Достоинства, недостатки и альтернативы объектно-реляционного отображения (ORM) / С.~С.~Романов.~--~Текст~: электронный // Таврический научный обозреватель.~--~2016.~--~№~11-2 (16).~--~URL: \url{https://cyberleninka.ru/article/n/dostoinstva-nedostatki-i-alternativy-obektno-relyatsionnogo-otobrazheniya-orm} (дата обращения: 10.05.2026).

\bibitem{phpdoc} PHP Documentation~: official manual.~--~The PHP Group, 1997--.~--~URL: \url{https://www.php.net/docs.php} (дата обращения: 10.05.2026).~--~Текст~: электронный.

\bibitem{laraveldoc} Laravel Documentation~: official site.~--~Taylor Otwell \& Laravel LLC, 2011--.~--~URL: \url{https://laravel.com/docs/12.x} (дата обращения: 10.05.2026).~--~Текст~: электронный.

\bibitem{sanctum} Laravel Sanctum~: API token authentication~// Laravel Documentation.~--~Taylor Otwell, 2020--.~--~URL: \url{https://laravel.com/docs/12.x/sanctum} (дата обращения: 10.05.2026).~--~Текст~: электронный.

\bibitem{reverb} Laravel Reverb~: official documentation~// Laravel Documentation.~--~Taylor Otwell, 2024--.~--~URL: \url{https://laravel.com/docs/12.x/reverb} (дата обращения: 10.05.2026).~--~Текст~: электронный.

% --- Svelte и фронтенд ---
\bibitem{sveltedoc} Svelte Documentation~: official site.~--~Rich Harris \& Svelte contributors, 2016--.~--~URL: \url{https://svelte.dev/docs} (дата обращения: 10.05.2026).~--~Текст~: электронный.

\bibitem{sveltekitdoc} SvelteKit Documentation~: official site.~--~Svelte contributors, 2020--.~--~URL: \url{https://kit.svelte.dev/docs} (дата обращения: 10.05.2026).~--~Текст~: электронный.

\bibitem{tailwind} Tailwind CSS Documentation~: official site.~--~Tailwind Labs, 2017--.~--~URL: \url{https://tailwindcss.com/docs} (дата обращения: 10.05.2026).~--~Текст~: электронный.

\bibitem{mdn} MDN Web Docs~: официальный сайт.~--~Mountain View~: Mozilla Foundation, 2005--.~--~URL: \url{https://developer.mozilla.org} (дата обращения: 10.05.2026).~--~Текст~: электронный.

\bibitem{flanagan} Флэнаган, Д. JavaScript. Полное руководство / Д.~Флэнаган~; перевод с английского.~--~7-е изд.~--~Москва~: ООО <<Диалектика>>, 2021.~--~720~с.~--~ISBN 978-5-907203-79-2.~--~Текст~: непосредственный.

\bibitem{robson} Робсон, Э. Изучаем HTML и CSS / Э.~Робсон, Э.~Фримен~; перевод с английского.~--~2-е изд.~--~Санкт-Петербург~: Питер, 2014.~--~720~с.~--~ISBN 978-5-496-00653-8.~--~Текст~: непосредственный.

% --- Программная инженерия, архитектура ---
\bibitem{orlov} Орлов, С.~А. Программная инженерия. Технологии разработки программного обеспечения~: учебник для вузов / С.~А.~Орлов.~--~5-е изд., обнов. и доп.~--~Санкт-Петербург~: Питер, 2020.~--~640~с.~--~ISBN 978-5-4461-1348-4.~--~Текст~: непосредственный.

\bibitem{fowler_uml} Фаулер, М. UML. Основы~: краткое руководство по стандартному языку объектного моделирования / М.~Фаулер~; перевод с английского.~--~3-е изд.~--~Санкт-Петербург~: Символ-Плюс, 2013.~--~192~с.~--~ISBN 5-93286-060-X.~--~Текст~: непосредственный.

\bibitem{fowler_enterprise} Фаулер, М. Шаблоны корпоративных приложений / М.~Фаулер, Д.~Райс, М.~Фоммел [и др.]~; перевод с английского.~--~Москва~: ООО <<И.~Д.~Вильямс>>, 2018.~--~544~с.~--~ISBN 978-5-8459-1611-2.~--~Текст~: непосредственный.

\bibitem{gamma} Гамма, Э. Паттерны объектно-ориентированного проектирования / Э.~Гамма, Р.~Хелм, Р.~Джонсон, Дж.~Влиссидес~; перевод с английского.~--~Санкт-Петербург~: Питер, 2020.~--~448~с.~--~(Библиотека программиста).~--~ISBN 978-5-4461-1595-2.~--~Текст~: непосредственный.

\bibitem{mikrotikwiki} MikroTik Documentation~: official wiki.~--~MikroTikls SIA, 2003--.~--~URL: \url{https://help.mikrotik.com/docs/} (дата обращения: 10.05.2026).~--~Текст~: электронный.

\bibitem{routerosapi} MikroTik RouterOS API~: protocol description // MikroTik Documentation.~--~MikroTikls SIA, 2007--.~--~URL: \url{https://help.mikrotik.com/docs/display/ROS/API} (дата обращения: 10.05.2026).~--~Текст~: электронный.

\bibitem{routerosphp} routeros-api-php~: PHP client for the MikroTik RouterOS API / P. Zhalybin (evilfreelancer).~--~GitHub, 2017--.~--~URL: \url{https://github.com/EvilFreelancer/routeros-api-php} (дата обращения: 10.05.2026).~--~Текст~: электронный.

\bibitem{burton} Burton, T. MikroTik RouterOS Workshop: QoS Best Practice / T.~Burton~; MikroTik User Meeting Presentation.~--~Riga~: MikroTikls, 2018.~--~URL: \url{https://mum.mikrotik.com/presentations/} (дата обращения: 10.05.2026).~--~Текст~: электронный.

\bibitem{olifer} Олифер, В.~Г. Компьютерные сети. Принципы, технологии, протоколы~: учебник для вузов / В.~Г.~Олифер, Н.~А.~Олифер.~--~5-е изд.~--~Санкт-Петербург~: Питер, 2016.~--~992~с.~--~ISBN 978-5-496-01967-5.~--~Текст~: непосредственный.

\bibitem{splynx} Splynx ISP Framework~: official website.~--~Splynx s.r.o., 2015--.~--~URL: \url{https://www.splynx.com/} (дата обращения: 10.05.2026).~--~Текст~: электронный.

\bibitem{sonar} Sonar Software~: official website.~--~Sonar Software, Inc., 2014--.~--~URL: \url{https://sonar.software/} (дата обращения: 10.05.2026).~--~Текст~: электронный.

\bibitem{uisp} UISP~: Ubiquiti ISP Management Platform.~--~Ubiquiti Inc., 2016--.~--~URL: \url{https://www.ui.com/uisp/} (дата обращения: 10.05.2026).~--~Текст~: электронный.

\bibitem{powercode} Powercode~: official website.~--~Powercode LLC, 2008--.~--~URL: \url{https://powercode.com/} (дата обращения: 10.05.2026).~--~Текст~: электронный.

\bibitem{ispconfig} ISPConfig~: hosting control panel.~--~projektfarm GmbH, 2005--.~--~URL: \url{https://www.ispconfig.org/} (дата обращения: 10.05.2026).~--~Текст~: электронный.

\bibitem{carbon} Carbon Billing 5~: официальный сайт.~--~ООО <<Карбон Софт>>, 2003--.~--~URL: \url{https://carbonsoft.ru/products/carbon-billing-5/} (дата обращения: 10.05.2026).~--~Текст~: электронный.

\bibitem{netup} NetUP UTM~5~: официальный сайт.~--~ООО <<НетАП>>, 2001--.~--~URL: \url{https://www.netup.ru/} (дата обращения: 10.05.2026).~--~Текст~: электронный.

% --- Информационная безопасность ---
\bibitem{uskov} Усков, И.~Д. Безопасность веб-приложений: XSS, CSRF и современные методы защиты / И.~Д.~Усков, В.~Д.~Сандаков.~--~Текст~: электронный // Вестник науки.~--~2026.~--~URL: \url{https://cyberleninka.ru/article/n/bezopasnost-veb-prilozheniy-xss-csrf-i-sovremennye-metody-zaschity} (дата обращения: 10.05.2026).

\end{thebibliography}
